\documentclass[a4paper,12pt]{article}

% Paquetes necesarios
\usepackage[utf8]{inputenc} % Codificación de caracteres UTF-8
\usepackage[spanish]{babel} % Soporte para español
\usepackage{graphicx} % Insertar imágenes
\usepackage{fancyhdr} % Encabezados y pies de página personalizados
\usepackage{tocbibind} % Para incluir la bibliografía en el índice
\usepackage{setspace} % Controlar el espacio entre líneas
\usepackage{enumitem} % Listas personalizadas
\usepackage{titlesec} % Formatear títulos
\usepackage{geometry} % Configurar márgenes
\usepackage{newtxtext,newtxmath} % Tipografía Times New Roman
\usepackage{ragged2e} % Alineación justificada
\usepackage{url} % Mostrar URLs
\usepackage{array} % Mejoras en las tablas
\usepackage{pgfgantt} % Diagramas de Gantt
\usepackage{lscape} % Páginas en modo apaisado
\usepackage{xcolor} % Colores personalizados
\usepackage{verbatim} % Entorno verbatim
\usepackage{listings} % Mostrar código fuente

% Configuración de márgenes
\geometry{
	a4paper,
	left=3.5cm,
	right=2.5cm,
	top=3.5cm,
	bottom=3.5cm
}

% Espaciado entre líneas
\onehalfspacing

% Configuración de los títulos para secciones
\titleformat{\section}
{\normalfont\bfseries\large} % Formato de sección: normal, negrita, 14pt (LARGE)
{\thesection}{1em}{\MakeUppercase} % Mayúsculas
\titlespacing*{\section}{0pt}{6pt}{18pt} % Espaciado de sección: 6pt antes, 18pt después
\newcommand{\sectionbreak}{\clearpage} % Nueva página para cada sección

% Configuración de los títulos para subsecciones
\titleformat{\subsection}
{\normalfont\bfseries\normalsize} % Formato de subsección: normal, negrita, 12pt (Large)
{\thesubsection}{1em}{}
\titlespacing*{\subsection}{0pt}{6pt}{18pt} % Espaciado de subsección: 6pt antes, 18pt después

% Configuración de los títulos para subsubsecciones
\titleformat{\subsubsection}
{\normalfont\bfseries\normalsize} % Formato de subsubsección: normal, negrita, 10pt (large)
{\thesubsubsection}{1em}{}
\titlespacing*{\subsubsection}{0pt}{6pt}{18pt} % Espaciado de subsubsección: 6pt antes, 18pt después

% Formato específico para Resumen y Abstract
\titleformat{name=\section,numberless}[block]
{\normalfont\bfseries\normalsize\centering}
{}
{0pt}
{}

% Capitalizar títulos
\addto\captionsspanish{
	\renewcommand{\contentsname}{Índice General}
	\renewcommand{\listfigurename}{Índice de Imágenes}
	\renewcommand{\listtablename}{Índice de Figuras}
	\renewcommand{\tablename}{Tabla}
	\renewcommand{\refname}{Referencias Bibliográficas}
}

% Configuración de los encabezados y pies de página
\pagestyle{fancy}
\fancyhf{}
\fancyhead{} % Eliminar contenido del encabezado
\fancyfoot[R]{\thepage} % Número de página en la esquina inferior derecha
\renewcommand{\headrulewidth}{0pt} % Eliminar la línea horizontal en la cabecera
\renewcommand{\footrulewidth}{0pt} % Eliminar la línea horizontal en el pie de página

% Redefinir numeración romana para minúsculas en las páginas preliminares
\fancypagestyle{plain}{
	\fancyhf{}
	\fancyfoot[R]{\thepage}
	\renewcommand{\thepage}{\roman{page}}
}

\begin{document}
	
	% Inclusión de la portada
	\include{secciones/portada}
	
	% Para numeración romana desde la Hoja de Aprobación hasta el Índice de Figuras
	\clearpage
	
	\pagenumbering{roman}    
	
	\setcounter{page}{2} % Iniciar en ii
	\renewcommand{\thepage}{\roman{page}}
	
	% Hoja de aprobación
	\pagestyle{plain}
	\include{secciones/hoja_aprobacion}
	\thispagestyle{plain} % Asegurarse de que la hoja de aprobación tenga número de página
	
	\newpage
	
	% Declaración de conformidad
	\pagestyle{plain}
	\include{secciones/declaracion_conformidad}
	\thispagestyle{plain} % Asegurarse de que la página tenga número de página
	
	% Índice General (excluido del TOC)
	\addtocontents{toc}{\protect\setcounter{tocdepth}{-1}}
	\tableofcontents
	\addtocontents{toc}{\protect\setcounter{tocdepth}{2}}
	\newpage
	
	% Índice de Tablas (excluido del TOC)
	\addtocontents{toc}{\protect\setcounter{tocdepth}{-1}}
	\listoftables
	\addtocontents{toc}{\protect\setcounter{tocdepth}{2}}
	\newpage
	
	% Índice de Figuras (excluido del TOC)
	\addtocontents{toc}{\protect\setcounter{tocdepth}{-1}}
	\listoffigures
	\addtocontents{toc}{\protect\setcounter{tocdepth}{2}}
	\newpage
	
	% Resumen o Abstract
	\pagestyle{plain}
	\include{secciones/resumen}
	\thispagestyle{plain}
	\newpage
	
	% Cambio a numeración arábiga
	\pagenumbering{arabic}
	\fancyfoot[R]{\thepage} % Número de página en la esquina inferior derecha
	
	% Justificar el contenido
	\justifying
	
	% Eliminar la sangria en los parrafos
	\setlength{\parindent}{0pt}
	
	% Espaciado entre párrafos de 6 puntos
	\setlength{\parskip}{6pt} 
	
	% Sección 1: Introducción
	\include{capitulos/introduccion}
	
	% Sección 2: Conceptos Fundamentales
	\include{capitulos/conceptos_fundamentales}
	
	% Sección 3: Materiales y Métodos
	\include{capitulos/materiales_y_metodos}
	
	% Sección 4: Análisis de Datos
	\include{capitulos/analisis_de_datos}
	
	% Sección 5: Aspectos que puedan anticiparse
	\include{capitulos/aspectos_anticipables}
	
	% Sección 6: Calendarización
	\include{capitulos/calendarizacion}
	
	% Referencias Bibliográficas
	\bibliographystyle{IEEEtran}
	\bibliography{bibliografia/referencias}
	
	% Sección 8: Anexos
	\addcontentsline{toc}{section}{Anexos} % Añadir Anexos al índice
	\include{capitulos/anexos}
	%\include{secciones/miscelaneas}
	
\end{document}
